\chapter{OBJETIVOS}

\section{Objetivo Geral}

O presente estudo tem como objetivo principal desenvolver, avaliar e comparar modelos de Aprendizado de Máquinas (AM) capazes de prever o desempenho de equipes,
a partir de dados estatísticos de desempenho individual de jogadores. Uma vez definido o modelo preditivo com melhor desempenho, este será utilizado como função
objetivo em modelos simples de otimização, com o intuito de auxiliar o processo de tomada de decisão em contratações, buscando identificar soluções de mercado
mais adequadas para a composição de elencos no futebol profissional.

\section{Objetivos Específicos}

Com o intuito de atingir o objetivo geral, foram estabelecidos os seguintes objetivos específicos:

\begin{enumerate}
    \item Coletar, estruturar e processar dados históricos de jogadores e partidas, abrangendo múltiplas temporadas, garantindo a qualidade e a integridade
    das informações utilizadas no desenvolvimento dos modelos preditivos.

    \item Desenvolver uma metodologia para quantificar o impacto individual dos jogadores, combinando estatísticas históricas ao longo das temporadas por meio
    de um sistema de ponderação temporal, priorizando desempenhos mais recentes.

    \item Construir modelos de AM para prever o desempenho agregado de uma equipe em uma determinada temporada, considerando a composição
    do elenco e os dados históricos dos jogadores.

    \item Validar a eficácia dos modelos preditivos, comparando as estimativas com os resultados reais de temporadas subsequentes, utilizando métricas
    estatísticas como erro médio absoluto (MAE) e coeficiente de determinação ($R^2$).

    \item Desenvolver modelos simples de otimização baseados no modelo de AM selecionado, visando a identificação de combinações de jogadores que maximizem
    o desempenho esperado da equipe sob restrições de mercado.
\end{enumerate}

Ao integrar técnicas de AM e otimização, este estudo busca fornecer uma ferramenta analítica de apoio à decisão para a montagem de
elencos no futebol profissional, contribuindo para práticas de gestão esportiva orientadas por dados.